\documentclass[12pt]{article}
\usepackage[utf8]{inputenc}
\usepackage[T1]{fontenc}
\usepackage{hyperref}
\title{Ancient North Arabian Language Resources for K-12 Curriculum Development}
\author{Compiled by Claude Sonnet 5.5 (claude-sonnet-5-5) for Andr\'as Kornai}
\date{}
\begin{document}
\maketitle

\section{Introduction}
Ancient North Arabian (ISO 639-3 xna) is a cover term for the languages written in the Ancient North Arabian scripts. These are Safaitic, Hismaic, Dadanitic, Taymanitic, Dumaitic and the Thamudic varieties. The scripts were used from roughly the early first millennium BC to perhaps the 5th century AD. The language has no living speakers or community. All material is epigraphic and scholarly. It consists of short inscriptions, and no native-language school materials exist.

\section{Primary Educational Resources}
I found no textbooks or primers. Teaching would have to rely on the corpus and scholarly works below.

\section{Linguistic and Reference Materials}
\begin{itemize}
\item OCIANA, the Online Corpus of the Inscriptions of Ancient North Arabia (Oxford). The project description gives its scope as all pre-Islamic inscriptions of north and central Arabia: \url{https://www.orinst.ox.ac.uk/node/142271}
\item OCIANA overview of the Ancient North Arabian scripts (Ohio State): \url{https://ociana.osu.edu/scripts_north_arabian}
\item Survey of the Ancient North Arabian inscriptions from the Dumat al-Jandal area, Saudi Arabia: \url{https://archaeopresspublishing.com/ojs/index.php/PSAS/article/view/1203}
\item Leiden University scholarly publication on Safaitic texts: \url{https://scholarlypublications.universiteitleiden.nl/access/item%3A3283748/view}
\item Ahmad Al-Jallad, pre-prints and open-access books, which include a 2025 handout on the Thamudic categories: \url{https://osu.academia.edu/AhmadAlJallad}
\end{itemize}

\section{Religious and Community Resources}
No community or religious resources apply. Jordan Times report on research into the Safaitic script: \url{https://www.jordantimes.com/news/local/new-insights-emerged-ancient-safaitic-script}

\section{Digital and Audio Resources}
The OCIANA database is free to use online and also offers preliminary PDF editions by script. No audio exists.

\section{Curriculum Development Context}
This is a language of history and epigraphy. A K-12 curriculum cannot be taught in it. A realistic unit would cover the scripts, the inscriptions and Arabian desert-nomad history, using OCIANA images and transcriptions. Scholars disagree on whether the scripts are genetically related, and the Thamudic material is the least uniform.
\end{document}